Deux moitiés que tout sépare. Les statistiques descriptives sont du calcul : la moyenne est
une somme divisée par un effectif, et ses propriétés — linéarité, effet d'une translation —
se démontrent en développant la somme. Les probabilités, elles, demandent un modèle.

C'est le premier chapitre où formaliser oblige à choisir. \og{} On tire une boule au hasard
\fg{} n'est pas une proposition mathématique tant qu'on n'a pas dit sur quel ensemble et avec
quelle loi. Le chapitre travaille donc avec des probabilités uniformes sur un ensemble fini,
ce qui est exactement le modèle du collège, et démontre ce qui en découle : la probabilité
d'un événement contraire, celle d'une réunion, l'équiprobabilité.
